% Generated draft: requires mathematical and source review.
\documentclass[11pt]{article}
\usepackage[margin=25mm]{geometry}
\usepackage{fontspec}
\setmainfont{Latin Modern Roman}
\usepackage{amsmath,amssymb,mathtools,unicode-math}
\setmathfont{Latin Modern Math}
\usepackage{enumitem,needspace,xurl,hyperref}
\setlength{\parindent}{0pt}
\setlength{\parskip}{5pt}
\setlength{\emergencystretch}{3em}
\newcommand{\N}{\mathbb{N}}
\newcommand{\Z}{\mathbb{Z}}
\newcommand{\Q}{\mathbb{Q}}
\newcommand{\R}{\mathbb{R}}
\newcommand{\C}{\mathbb{C}}
\newcommand{\F}{\mathbb{F}}
\newcommand{\A}{\mathbb{A}}
\newcommand{\PP}{\mathbb{P}}
\newcommand{\E}{\mathbb{E}}
\newcommand{\eps}{\varepsilon}
\newcommand{\dd}{\,\mathrm{d}}
\begin{document}
\section*{UnsolvedMath ID: 30000437}
\textbf{Source problem code:} catalogue entry
\subsection{Definitions and mathematical statement}
The source asks whether one can cut the boundary of a convex polytope along edges and begin an infinitesimal unfolding motion during which every face remains on the convex hull. The natural universal interpretation is: for every admissible convex polytope, does some edge-cut choice admit such a motion? The dimension, admissible cut patterns and precise meaning of infinitesimal are not stated.

For a concrete conventional notation in dimension three, let \(P\subset\R^3\) be a full-dimensional convex polytope with facets \(\mathcal F\), and choose a subset of its original edges to cut. An actual local rigid-face motion would assign each facet a family of Euclidean rigid motions \(T_F(t)\), for \(0\leq t<\delta\), starting from \(T_F(0)=\operatorname{id}\). Along every uncut common edge the corresponding copies of its points must remain identified, so the edge serves as a hinge. Put
\[
 U_t=\bigcup_{F\in\mathcal F}T_F(t)F,\qquad K_t=\operatorname{conv}(U_t).
\]
A natural full-face interpretation of the hull requirement is
\[
 T_F(t)F\subseteq\partial K_t\quad
 \text{for every facet }F\text{ and every }0\leq t<\delta.
\]
The convex hull here is that of the moving faces, not the original solid \(P\). This is a notation scheme, conditional on the intended three-dimensional rigid-facet setting. Requiring all points of a face to remain on the boundary is stronger than requiring only its vertices to lie on that boundary.

An unfolding must involve genuine relative motion rather than a common ambient rigid motion. But the source does not say whether it seeks an actual nontrivial differentiable path on a positive time interval or only first-order velocities satisfying linearized constraints. First-order feasibility does not by itself supply an integrable motion. The permitted hinge-angle changes and any monotone-opening requirement are not supplied. Neither are connectedness of the cut surface, a spanning-tree cut condition, or retention of a single connected hinged assembly; allowing all facets to be separated would materially change the question.

The requested initial motion is not automatically a continuous unfolding all the way to a planar nonoverlapping net. The cleaned/original sentence does not explicitly impose such a final state or a collision convention. Keeping every face on a convex hull also does not by itself rule out coincident or overlapping boundary facets. The 2026-08-21 literature triage discusses the broader nonoverlapping blooming problem using only original edges, reporting known source-unfolding and refinement results but an unresolved edge-only problem. Those results and stronger global conditions are kept as context, not treated as equivalent to the local hull-preserving request. The exact-statement verification and open label remain source metadata.
\subsection{Short English statement}
Can every intended convex polytope be cut along original edges so that a genuine initial unfolding keeps every moving face on the boundary of their convex hull? The source leaves the cut connectivity and infinitesimal-motion definition unspecified. Its local question is distinguished from the stronger global nonoverlapping blooming problem discussed in the literature note.
\subsection{Sources}
\begin{itemize}
\item Local UnsolvedMath registry, permanent ID 30000437, source code catalogue entry, ``Convex-Hull-Preserving Bloomings of Polytope Boundaries,'' including identical original and cleaned local-motion question and verification metadata.
\item \emph{Discrete Differential Geometry (2006)}, Oberwolfach Reports, pp.\ 653--728, DOI 10.4171/owr/2006/12: \url{https://doi.org/10.4171/owr/2006/12}.
\item Supplied reference ``Demaine et al. continuous blooming'': \url{https://doi.org/10.48550/arXiv.0906.2461}.
\item Literature triage checked 2026-08-21; registry record updated 2026-09-16. No external source has been visited or independently checked here.
\end{itemize}
\end{document}
